\title{LongRoPE: Extending LLM Context Window Beyond 2 Million Tokens}

\begin{abstract}
A large context window remains a valuable characteristic in large language models (LLMs). Nonetheless, existing approaches are constrained to approximately 128k tokens because of the considerable expense of fine-tuning, limited availability of lengthy textual data, and the emergence of problematic values associated with novel token positions.
\end{abstract}

This study presents LongRoPE, marking the first time that pre-trained LLMs have achieved an extensive context window of \textbf{2048k} tokens—accomplished with as few as 1k fine-tuning iterations at training lengths of up to 256k, all while preserving their baseline performance at the original shorter context window. This advancement is underpinned by three primary contributions: (i) we discover and leverage two types of non-uniformities within positional interpolation, identified through an efficient search process, thereby facilitating enhanced initialization for fine-tuning and enabling an 8$\times$ context length expansion even in settings without fine-tuning; (ii) we propose a staged extension methodology where a model is first fine-tuned at 256k length, followed by a secondary positional interpolation on the extended, fine-tuned model, ultimately yielding a 2048k context window; (iii) we refine LongRoPE at 8k length to restore original short context window effectiveness. Comprehensive experiments on both LLaMA2 and Mistral, covering a diverse array of tasks, validate the efficacy of our approach. Importantly, models adapted with LongRoPE preserve the initial architecture with only minimal changes limited to the positional embeddings, allowing the continued use of most existing optimizations. The code will be released at \texttt{https://github.com/microsoft/LongRoPE}

\section{Introduction}

Large Language Models (LLMs) have achieved impressive results across a range of tasks~\cite{gpt4,llama2}; however, they are constrained by their finite context window size, such as the 4096-token maximum of LLaMA2~\cite{llama2}. When input extends beyond this context window, LLM performance deteriorates because the model has not been trained to handle additional sequence positions. This limitation creates difficulties in key applications, including in-context learning where many examples are provided~\cite{Huang2023FewerIM}, as well as in LLM-driven agents~\cite{park2023generative,madaan2023selfrefine}.

Recent studies have demonstrated that it is possible to expand the context window of a pre-trained LLM to approximately 128k tokens by further training the model on longer textual sequences~\cite{longlora,pi,yarn,zhang2024soaring,liu2023scaling}. Nevertheless, three main challenges hinder additional increases in context window length. To begin with, integrating previously unused positional indices introduces a multitude of anomalous values, which gives rise to distributional shift problems and hampers the fine-tuning process, often making convergence difficult~\cite{pi}. This issue is especially pronounced when extending the context window from 4k to $>$1000k, since over 90\% of the token positions involved are novel. Second, effective fine-tuning generally necessitates access to training texts that match the targeted sequence length. Such ultra-long documents, particularly those with more than 1000k tokens, are scarce in current datasets. Additionally, training on extremely long sequences dramatically increases computational demands, making the process resource-intensive both in terms of training duration and required GPU capacity. Third, as the context window scales to much greater lengths, the model’s attention distribution becomes diluted across many token positions, which results in diminished performance on tasks requiring the shorter, original context length~\cite{pi}.

A commonly used strategy to address the first challenge is to apply interpolation to RoPE positional embeddings~\cite{rope,pi}, which rescales novel position indices to fit within the pretrained domain, as illustrated in Fig.~\ref{fig:overview}. The Position Interpolation (PI) method~\cite{pi} performs linear interpolation of the rotary angles in RoPE, determined by the scaling ratio. NTK~\cite{ntk,dynamicntk} introduces non-uniform interpolation and extrapolation along the dimensions of RoPE. In YaRN~\cite{yarn}, RoPE dimensions are segmented into three groups according to frequency, with extrapolation, NTK, and linear interpolation applied to each group, respectively. Nevertheless, positional embeddings possess a \emph{highly intricate and non-uniform information entropy} within the Transformer architecture. The nuanced nature of this non-uniformity is largely overlooked by current methods, resulting in the loss of valuable information and imposing constraints on the maximum extendable context window.

Section~\ref{sec:analysis} demonstrates through empirical evidence two principal observations: \textbf{(1)} Successful positional interpolation necessitates accounting for two distinct types of non-uniformities, specifically fluctuations across RoPE dimensions and inconsistencies in token positions. Reduced RoPE dimensions and early sequence token locations generally require less aggressive interpolation, yet the most suitable approach is influenced by the particular extended length aimed for.  
\textbf{(2)} Integrating awareness of these non-uniformities into positional interpolation techniques helps to better preserve the informational content embedded within the original RoPE, especially in terms of critical dimensions and token locations. This approach reduces the information loss typically induced by positional interpolation and consequently offers a more effective initialization prior to fine-tuning. Additionally, it supports an 8$\times$ sequence extension even in the absence of fine-tuning.

Building on these results, we propose \textbf{LongRoPE}, a robust approach for pushing the context window of LLMs past 2 \emph{million} tokens. The foundation of LongRoPE rests on three principal contributions. Initially, {LongRoPE} takes full advantage of the multidimensional, non-uniform characteristics inherent to positional interpolation. This is achieved by discovering optimized rescale factors for the rotational angles in RoPE along each dimension, guided by the distribution of token positions. Given that the space for potential rescale factors grows exponentially relative to the desired extension scale, {LongRoPE} employs an evolutionary search algorithm enhanced with two specialized optimization strategies to significantly improve search efficiency. An illustrative instance of the optimized, rescaled RoPE discovered by this process is displayed in Fig.~\ref{fig:overview}.

Subsequently, {LongRoPE} adopts a resource-efficient and incremental extension approach to realize a 2048k context window, circumventing the necessity for direct fine-tuning on ultra-long text data, which is typically rare and challenging to obtain. Initially, this process involves identifying a 256k sequence length using the pre-trained LLM and performing fine-tuning at this scale. Given that our non-uniform positional interpolation method supports an 8$\times$ length extension without additional fine-tuning, we proceed to perform a second search for updated RoPE rescaling factors using the already fine-tuned, length-extended LLM. As a result, a 2048k context window is achieved for LLaMA2 and Mistral~\cite{mistral}.

Lastly, to address potential drops in performance on the initial (shorter) context window, {LongRoPE} further refines the RoPE rescaling factors within the extended LLM. Following the same approach used for expanding from 256k to 2048k, we apply our search method to downscale the context window to 4k and 8k on the model fine-tuned for 256k, aiming to reduce excessive positional interpolation. At inference time, when dealing with sequence lengths below 8k, we replace the RoPE parameters with the optimized rescale factors identified in the search.

Comprehensive experimentation on multiple LLMs and a range of long-context tasks validates the efficacy of our approach. Our results indicate that LongRoPE consistently sustains low perplexity across evaluation lengths from 4k up to 2048k, attains more than 90\% accuracy in passkey retrieval, and matches the performance of baseline models on established benchmarks that utilize a 4096-token context window. Furthermore, LongRoPE is compatible with any LLM utilizing RoPE embeddings. We intend to publicly release both our code and the LongRoPE-2048k models.

\section{Non-uniformity in Positional Interpolation}
\label{sec:analysis}

\subsection{Preliminary}

Transformer models require explicit positional information, often in the form of position embedding, to represent the order of input tokens. Our work focuses on the RoPE~\cite{rope} position embedding, which is widely used in recent  LLMs. For a token at position index $n$, its corresponding RoPE encoding can be simplified as follows:

\begin{equation}
		\fontsize{7.8}{7.8} \selectfont
	\label{eq:rope}
	\left[ cos(n\theta_0), sin(n\theta_0), cos(n\theta_1), \cdot\cdot\cdot, cos(n\theta_{d/2-1}), sin(n\theta_{d/2-1}) \right]   
\end{equation}

where $d$ denotes the dimension of the embeddings, 
$n\theta_i$ refers to the rotary phase corresponding to the token situated at position $n$, 
while $\theta_i=\theta^{-2i/d}$ encodes the rotation frequency. The typical choice for the base $\theta$ in RoPE is 10000 by default.

\textbf{\emph{Context window extension ratio $s$} and positional interpolation}. We define the extension ratio $s$ as the quotient of the extended context length $L'$ over the original context length $L$, given by $s=\frac{L'}{L}$.

To extend the context window from $L$ to $L'$, current positional interpolation methods suggest downscaling rotation frequencies $\theta_i$ based on extension ratio $s$. Let $\beta=\theta^{2/d}$, and $\lambda$ denote the actual rescale factor related to $s$, we   unify these positional interpolation methods as follows:

\begin{equation}	
	\fontsize{7.5}{7.5} \selectfont
	\label{eq:unified_rope}
	\begin{aligned}
		\left[cos\left(\frac{n}{\lambda(\beta)^0}\right), sin \left(\frac{n}{\lambda(\beta)^0}\right), cos\left(\frac{n}{\lambda(\beta)^1}\right), \cdot\cdot\cdot,  sin \left(\frac{n}{\lambda(\beta)^{d/2-1}}\right) \right]   
	\end{aligned}
\end{equation}

\textbf{Linear positional interpolation (PI)}. In PI~\cite{pi}, position indices are interpolated linearly within the maximum sequence length seen during pre-training. To achieve an extension by a factor of $s$, the method applies a uniform scaling of the rotation angles in all RoPE dimensions by a factor $\lambda=s$. Despite its simplicity, this approach leads to a compression of positional representations, causing tokens with nearby positions to be less distinguishable. As a result, PI often yields suboptimal results when used with larger extension ratios.

\textbf{NTK-based interpolation and extrapolation}. Works such as~\cite{ntk,dynamicntk} approach RoPE through the lens of information encoding, leveraging the Neural Tangent Kernel (NTK) framework~\cite{jacot2018neural,tancik2020fourier}. In response to the problem of overcrowding in position encoding (PI), these methods advocate distributing interpolation demands unevenly across RoPE’s different dimensions. Specifically, dimensions associated with higher frequencies (lower indices) are scaled down less, while those with lower frequencies (higher indices) are scaled up more, enabling both interpolation and extrapolation over positional encodings, described mathematically as $\lambda=s^{i}$. The refined dynamic NTK approach~\cite{dynamicntk} further tunes the scaling factor depending on the current sequence’s length at each position. In contrast to PI, which requires subsequent fine-tuning, NTK-based strategies permit context window expansion without fine-tuning; however, they generally limit the extension ratio to 4$\times$.

\textbf{YaRN}~\cite{yarn} presents a notable advance in the effectiveness of positional interpolation. The method segments the RoPE dimensions into three separate groups according to their associated frequencies, assigning a unique interpolation technique to each group. Extrapolation ($\lambda$=1) is reserved for the highest frequency dimensions, whereas linear interpolation (PI) is applied to dimensions with the lowest frequencies. For the dimensions lying between these extremes, the NTK approach is used. A central feature of YaRN is its strategy for categorizing RoPE dimensions, which relies primarily on heuristic, manual experimentation. This reliance on empirical grouping could hinder optimal performance when applied to newly-developed LLM architectures.

\subsection{Study on Non-uniform Positional Interpolation}

Drawing on insights from NTK and YaRN, we observe that their improvements stem from leveraging non-linear behaviors, particularly by treating various frequency components along RoPE axes separately to enable targeted interpolation and extrapolation. Nonetheless, prevalent approaches to non-linearity are predominantly governed by manually crafted heuristics. This observation prompts two main inquiries: (1) Does the existing approach to positional interpolation represent the best possible solution? (2) Might there exist additional, yet unidentified, forms of non-linearity?

To investigate these questions, we employ evolution search (refer to Sec.~\ref{sec:method}) to identify more effective non-uniform positional interpolations for LLaMA2-7B. Perplexity serves as the optimization criterion, measured over 5 randomly selected samples from the PG19~\cite{pg19} validation set. Our experiments yield the following significant insights.

\textbf{Finding 1}: \textit{RoPE dimensions exhibit substantial non-uniformities, which are not effectively handled by current positional interpolation methods.}

We determine the best value of $\lambda$ for each RoPE dimension as specified in Eq.~\ref{eq:unified_rope}. In Table~\ref{tbl:exp1}, we present a comparison of LLaMA2-7B's perplexity across various methods on the PG19 and Proof-pile~\cite{proof-pile} test sets, all evaluated without any fine-tuning. Our optimized approach yields notable perplexity reductions, indicating that commonly used linear (PI) and non-uniform (Dynamic-NTK and YaRN) interpolation strategies are less effective. Particularly, YaRN falls short relative to PI and NTK on the PG19 dataset, failing to achieve the intended context window for non-fine-tuned language models; for instance, with YaRN, perplexity increases sharply after 7k tokens in an 8k-length context.

In our exploration, the rescaling factors $\lambda$ in Eq.~\ref{eq:unified_rope} emerge as non-uniform, in contrast to the constant scale $s$ utilized in PI, the computation in NTK, and the group-specific scaling in YaRN. Adopting these varied scaling factors leads to a notable enhancement in LLaMA2's ability to model language—measured by perplexity—across 8k and 16k context lengths, even in the absence of any further fine-tuning. This improvement stems from the modified positional embeddings' ability to better maintain the structure of the original RoPE, particularly in crucial dimensions, thereby making it easier for the model to distinguish between tokens located close together.

\textbf{Finding 2}: \textit{RoPE for the initial tokens in the input sequence should be extrapolated with less interpolation.} 

We propose that for the first $\hat{n}$ tokens in the input sequence, the RoPE interpolation should be minimized. This is due to their disproportionately high attention scores, which renders them particularly significant within attention mechanisms—a pattern also reported in Streaming LLM~\cite{streamingllm} and LM-Infinite~\cite{han2023lminfinite}. To assess this hypothesis, we conduct experiments where the context window is expanded to 8k and 16k utilizing PI and NTK methods, while selectively avoiding interpolation for the initial $\hat n$ tokens ($\hat n$ = 0, 2, ..., 256). When $\hat n$ equals 0, the method is equivalent to standard PI and NTK. The results in Table~\ref{tbl:tokenposition} yield two major insights: \textbf{(1)} Exempting the earliest tokens from position interpolation enhances overall performance. \textbf{(2)} The ideal value of $\hat n$ varies as a function of the target context window size.

\textbf{Finding 3}: \textit{Non-uniform positional interpolation effectively extends LLM context window in both fine-tuning and non-fine-tuning settings.} 

Although our experiments demonstrate that the proposed non-uniform position interpolation markedly enhances extension performance at 8k and 16k context lengths even without fine-tuning, extending to even longer contexts still necessitates fine-tuning. Therefore, we conducted fine-tuning on LLaMA2-7B using our discovered RoPE configuration for a 64k context window (refer to the Appendix for specific hyperparameters). As reported in Table~\ref{tbl:finetune}, our approach consistently outperforms both PI and YaRN, for both the base and fine-tuned LLaMA2-7B models. This advantage is attributed to our approach’s more efficient non-uniform positional interpolation, which reduces information loss and yields superior parameter initialization for subsequent fine-tuning.

\textbf{Summary}. In our work, we identify and leverage two key non-uniformities: variations in RoPE dimensions and in token positions. Effectively exploiting these aspects in positional interpolation leads to substantial gains in LLM context extension capabilities.

\section{LongRoPE}

\label{sec:method}
Building on these insights, we propose LongRoPE. This method initially incorporates an effective search algorithm designed to leverage both forms of non-uniformity. Utilizing this approach, LongRoPE successfully extends the context window of LLMs to over 2 million tokens.

\subsection{Problem Formulation}

These dual non-uniformities contribute to an expansive solution landscape, thereby complicating the optimization process. To tackle this challenge, we recast the multidimensional non-uniform position interpolation optimization issue in terms of a search problem.

For a LLM targeting a  context window size of $L'$ and lengthy input documents $\mathbf{X}$, where each $\mathbf{x}\in\mathbf{X}$ surpasses $L'$ in token length, we denote the original rotary angle of the $i^{th}$ dimension in RoPE embedding at token position $n$ as  $\frac{n}{\beta^i}$. The optimization problem is then formulated as follows:

\begin{equation}
\begin{gathered}
		\label{eq:problem}

\underset {\mathbf{x}\in\mathbf{X}; \, |\mathbf{x}|\ge L'}{ \operatorname {arg\,min} } \,  \mathcal{L}\, \left( \text{LLM} (\text{RoPE}, \mathbf{X})\right),	\text{where} 
\\
\scalebox{0.93}{
$\underset {\begin{subarray}{l} i=0,\cdot\cdot,\frac{d}{2}-1;\\ n\in[ 0,|\mathbf{x}|);\end{subarray}}{\text{RoPE}_(n)}= {\left[\cdot\cdot,cos\left(\mathbb{I}(\hat{\lambda}_{i}, \hat{n})\times\frac{n}{\beta^i}\right),  sin\left(\mathbb{I}(\hat{\lambda}_{i}, \hat{n})\times\frac{n}{\beta^i}\right),\cdot\cdot\right] }$}\\
	\text{with} \quad \mathbb{I}(\hat{\lambda}_{i},\hat{n}) = \begin{cases}
	1&\text{if}\;  n<\hat{n}\\
	\frac{1}{\lambda_{i}}&\text{if}\; n\geq\hat{n}
\end{cases}
	\end{gathered}
\end{equation}

Here, we introduce a set of rescale factors, $\mathbb{I}(\hat{\lambda}_{i},\hat{n})$, designed to address both types of non-uniformity. The parameters $\hat{\lambda}_{i}$ and $\hat{n}$ correspond to non-uniformity along the RoPE dimensions and across token positions, respectively. More precisely, $\mathbb{I}(\hat{\lambda}_{i},\hat{n})$ serves to adjust the rotation angle for the $i^{th}$ RoPE dimension, where $\hat\lambda_{i}$ determines the scaling and $\hat{n}$ acts as the position cutoff. For token positions prior to $\hat{n}$ (i.e., for the first $\hat{n}$-1 tokens), no scaling is applied and the standard RoPE rotation angle $\frac{n}{\beta^i}$ is used. When $n\ge\hat{n}$, the rescale factor $\hat\lambda_{i}$ is enacted.

Assuming a desired context window length of $L'$, our goal is to identify the best set of rescale factors—specifically, ($\mathbb{I}(\hat{\lambda}_{0},\hat{n})$, $\mathbb{I}(\hat{\lambda}_{1},\hat{n})$, ..., $\mathbb{I}(\hat{\lambda}_{i},\hat{n})$, ...)—corresponding to each of the $d$ dimensions of the RoPE. The intent is that, once the $\text{RoPE}$ in the target $\text{LLM}$ has been adjusted accordingly, it will minimize the next-token prediction loss $\mathcal{L}$ (or, equivalently, the perplexity) for input sequences $\mathbf{X}$ consisting of $L'$ tokens.

\subsection{Searching the Non-uniform Position Interpolation}

To address the issue outlined in Eq.~\ref{eq:problem}, we present a straightforward but powerful approach that identifies the most suitable RoPE rescale factors, enabling comprehensive utilization of the multidimensional heterogeneity inherent in position embeddings.

\textbf{Search space}. We construct an extensive search space to accommodate the two forms of non-uniformity. The composition of this search space is depicted in Table~\ref{tbl:searchspace}. More precisely, our approach enables the optimization of individualized rescale factors for each separate dimension within the RoPE embedding. For ease in configuring the search space, we opt to search over $\lambda_i$ and $\hat{n}$ rather than directly tuning $\mathbb{I}(\hat{\lambda}_{i},\hat{n})$, where $\hat{\lambda}_{i}=1/\lambda_i$. According to Table~\ref{tbl:searchspace}, the search for $\lambda_i$ spans a range beginning at 1.0 (representing straightforward extrapolation) and extends up to $s\times1.25$ (corresponding to interpolation beyond what PI employs), incremented by 0.01; here, $s$ denotes the factor by which the context window is to be expanded.

The parameter $\hat{n}$ determines how many of the earliest token positions are kept with their original RoPE embeddings, bypassing position interpolation. In practice, we let $\hat{n}$ range over \{0, 1, 2, 4, 8, 12, 16, 20, 24, 28, 32, 64, 128, 256\}. Setting $\hat{n}=0$ means that rescale factors found through the search are applied to every token position.

\textbf{Evolution-based search}. The search space described in Table~\ref{tbl:searchspace} encompasses a vast array of positional interpolation possibilities, making it particularly difficult to navigate effectively. For instance, when the extension ratio is set to $s=4\times$, the number of possible configurations reaches $400^{128/2}\times14$ = 4$\times10^{167}$, and it grows exponentially as the extension ratio increases. To efficiently traverse this immense search space, we employ evolution search~\cite{guo2020single} and incorporate two key optimization strategies that substantially improve the efficiency of the search process. The entirety of our approach is detailed in Algorithm~\ref{alg:search}.

\textit{Optimized initial population generation}. Rather than generating all $P$ rescale factors in the initial population through random selection, we explicitly include the three RoPE rescale factors associated with PI, NTK, and YaRN as part of the initial pool. The other $P-3$ individuals are created by applying random mutations to these three rescale factors, with each factor mutated independently at probability $p$.

\textit{Monotonically non-decreasing constraint}. Following the creation of the initial population, we evaluate LLM perplexity for every candidate. In this step, we integrate the designated RoPE rescale factors into the target LLM and then determine the perplexity associated with the input $\mathbf{X}$. The highest-performing $k$ individuals are selected as parents for the subsequent evolutionary steps. Nonetheless, due to the extensive search space, unrefined mutation and crossover operators can often result in the exploration of suboptimal regions, requiring numerous unnecessary perplexity evaluations. This issue becomes more pronounced for large $L'$, where each perplexity computation entails a significant computational overhead.

To tackle this issue, we enforce a non-decreasing monotonicity condition on the sampled RoPE rescaled factors, namely, $\lambda_i\le \lambda_{i+1}$. Only those RoPE configurations meeting this requirement are considered for perplexity measurement in LLMs, which substantially streamlines the search process. More precisely, our approach dictates that $\lambda_i$ should monotonically increase with respect to the RoPE dimension index $i$ (where $i=0,\ldots,63$). This monotonicity along the dimension axis is inspired by insights from NTK theory~\cite{jacot2018neural,tancik2020fourier,ntk}. Specifically, the theory indicates that the lower-indexed, higher-frequency dimensions benefit from smaller $\lambda_i$ values due to a reduced need for interpolation, whereas the higher-indexed, lower-frequency dimensions can tolerate, or benefit from, increased interpolation using larger $\lambda_i$ values.

\begin{algorithm}[t]
	\small
	\caption{The search algorithm}
	\textbf{Input:} target LLM, input samples $\mathbf{X}$,   population size $P$, mutation size $N_1$, crossover size $N_2$, max iterations $\mathcal{T}$, mutate probability $p$\\

	\begin{algorithmic}[1]
		\label{alg:search}
		\STATE Top-k $\leftarrow \phi$;	
		\STATE $\text{P}_0$ $\leftarrow$ \textit{Initialize\_population\_with\_optimization}($P$, $\mathbf{X}$, $p$); 
		\FOR{$i = 1$ to $\mathcal{T}$}
		\STATE \textit{Compute\_perplexity}(LLM, $\text{P}_{i-1}$, $\mathbf{X}$);
		\STATE Top-k $\leftarrow$ \textit{Update\_Topk}(Top-k, $\text{P}_{i-1}$);
		\STATE $\text{P}_{mutation}$ $\leftarrow$ \textit{Mutation\_with\_mono\_constraint}(Top-k, $N_1$, $p$); 
		\STATE $\text{P}_{crossover}$ $\leftarrow$ \textit{Crossover\_with\_mono\_constraint}(Top-k, $N_2$);
		\STATE $\text{P}_i$ $\leftarrow$ $\text{P}_{mutation} \cup \text{P}_{crossover} \cup$ Top-k;
		\ENDFOR
		\STATE Return the member in Top-k achieving minimal perplexity;
	\end{algorithmic}
\end{algorithm}

\textbf{8$\times$ extension without fine-tuning}. Through evolutionary search, we successfully determine non-uniform RoPE rescale factors that protect critical dimensions and positions, thereby reducing information loss from interpolation. As shown in Fig.\ref{fig:non-ft}, our approach enables LLaMA2’s context window to be scaled from 4k up to 32k without the need for fine-tuning. By comparison, current techniques like PI, as well as non-uniform NTK and YaRN, result in significant increases in perplexity after merely a 2$\times$ context window extension.

\subsection{Extending LLM Context Window to 2048K}
\label{sec:extendto2048k}

\textbf{Progressive extension to 2048k}. In this section, we present our approach for scaling the context window of pre-trained LLMs far beyond the standard 4k, reaching lengths exceeding 2048k. Our non-uniform positional interpolation technique enables an 8$\times$ context extension with no need for fine-tuning. However, when scaling up to much higher factors (such as 512$\times$), fine-tuning becomes indispensable. One possible strategy involves searching for suitable RoPE rescaling parameters for the target window of 2048k tokens, followed by fine-tuning the model. This approach, however, encounters substantial obstacles due to the extremely high computational and resource demands required for such large-scale training. Furthermore, as reflected in our experimental observations, achieving robust fine-tuning results at these large extension ratios presents significant difficulties (refer to Appendix).

Fortunately, LongRoPE demonstrates efficacy on both base and fine-tuned extended LLMs. To this end, we propose a streamlined, incremental approach that reaches the desired 2048k context window using only 1k fine-tuning iterations, and operates within a maximum training length of 256k.

$\Diamond$ \textit{Expanding pre-trained LLMs to 256k using LongRoPE search}. Using LLaMA2 as a representative model, we perform a search procedure targeting context window sizes of 128k and 256k. In this process, the context extension ratios are 32$\times$ and 64$\times$, correspondingly.

$\Diamond$ \textit{Fine-tuning to 256k}. In this stage, we adapt the pre-trained LLM to support a context window of 256k tokens. To do so, we initially perform 400 fine-tuning steps on LLaMA2 with the RoPE rescaled factors adjusted for 128k. After these steps, we update the RoPE rescaled factors to correspond to 256k on the resulting checkpoint, followed by a further 600 fine-tuning steps. This approach is observed to be more efficient compared to commencing the fine-tuning process directly at 256k.

$\Diamond$ \textit{Scaling a fine-tuned extended LLM to 2048k via LongRoPE search}. In the final stage, we apply a subsequent search procedure to the LLM that has already been fine-tuned to handle 256k-length sequences. This approach enables us to reach a massive context window of 2048k, all without the need for any additional fine-tuning. The resulting context window achieves an expansion factor of $512\times$.

\textbf{Performance within the original context window}. When expanding the context window to an extensive 2048k length, we observe a degradation in performance on the model’s initial context range. This decline is a documented challenge associated with positional interpolation~\cite{pi}, which compels the position embeddings corresponding to the original context window to cluster in a significantly smaller region of the embedding space, thus impairing the language model’s effectiveness. With an extension factor of 512$\times$, positions spanning the original 4k context window become highly congested.

To address this issue, an additional evolution search is conducted on the expanded LLM, specifically tuning the RoPE rescale factors for shorter context lengths such as 4k and 8k. The upper bound for the search parameter $\lambda$ is decreased because shorter sequences necessitate minimal positional interpolation. While inferring, the LLM adapts the appropriate RoPE rescale factors in real time.

\section{LongRoPE}

\label{sec:method}
Building upon these observations, we propose LongRoPE, which initially incorporates a search algorithm designed to efficiently leverage the two identified non-uniformities. Utilizing this approach, we are able to extend the LLM's context window to exceed 2 million tokens.

\subsection{Problem Formulation}

These two forms of non-uniformity can generate an extensive set of possible solutions, thereby complicating the optimization process. To mitigate this, we reformulate the multidimensional non-uniform position interpolation optimization as a search problem.

For a LLM targeting a  context window size of $L'$ and lengthy input documents $\mathbf{X}$, where each $\mathbf{x}\in\mathbf{X}$ surpasses $L'$ in token length, we denote the original rotary angle of the $i^{th}$ dimension in RoPE embedding at token position $n$ as  $\frac{n}{\beta^i}$. The optimization problem is then formulated as follows:

\begin{equation}
\begin{gathered}
		\label{eq:problem}

\underset {\mathbf{x}\in\mathbf{X}; \, |\mathbf{x}|\ge L'}{ \operatorname {arg\,min} } \,  \mathcal{L}\, \left( \text{LLM} (\text{RoPE}, \mathbf{X})\right), \ \text{with}
\\
\scalebox{0.93}{
$\underset {\begin{subarray}{l} i=0,\ldots,\frac{d}{2}-1;\\ n\in[ 0,|\mathbf{x}|);\end{subarray}}{\text{RoPE}_{(n)}}= \left[\ldots,\cos\left(\mathbb{I}(\hat{\lambda}_{i}, \hat{n})\cdot\frac{n}{\beta^i}\right), \sin\left(\mathbb{I}(\hat{\lambda}_{i}, \hat{n})\cdot\frac{n}{\beta^i}\right),\ldots\right] $}\\
	\text{where \;} \mathbb{I}(\hat{\lambda}_{i},\hat{n})=
\begin{cases}
	1&\text{if } n<\hat{n}\\
	{\frac{1}{\lambda}_{i}}&\text{if } n\geq\hat{n}
\end{cases}
	\end{gathered}
\end{equation}

In this context, we define a group of scaling coefficients, $\mathbb{I}(\hat{\lambda}_{i},\hat{n})$, to address two types of non-uniformities. The parameters $\hat{\lambda_i}$ and $\hat{n}$ represent the dimensional non-uniformity in RoPE and the non-uniform distribution across token positions, respectively. More precisely, $\mathbb{I}(\hat{\lambda}_{i},\hat{n})$ is applied to adjust the rotation angle for the $i^{th}$ dimension of RoPE, utilizing $\hat{\lambda}_{i}$ as the scaling parameter and $\hat{n}$ as the threshold for token position. For token positions preceding $\hat{n}$, i.e., for the first $\hat{n}-1$ tokens, the rescaling coefficient $\hat{\lambda}_{i}$ remains inactive, maintaining the standard RoPE rotation angle $\frac{n}{\beta^i}$. The scaling factor is engaged only for tokens located at positions $n \ge \hat{n}$.

Assuming a desired context window length of $L'$, the task is to determine the set of optimal rescaling coefficients ($\mathbb{I}(\hat{\lambda}_{0},\hat{n})$, $\mathbb{I}(\hat{\lambda}_{1},\hat{n})$, ... $\mathbb{I}(\hat{\lambda}_{i},\hat{n})$, ...) corresponding to each RoPE dimension from $1$ through $d$. This rescaling enables the target $\text{LLM}$, utilizing the modified $\text{RoPE}$, to attain the lowest possible next-token prediction loss $\mathcal{L}$—equivalently, the lowest perplexity—across input samples $\mathbf{X}$ of token length $L'$.

\subsection{Searching the Non-uniform Position Interpolation}

In order to address the challenge outlined in Eq.~\ref{eq:problem}, we present our straightforward but powerful approach. This technique identifies the best RoPE rescaling parameters, enabling comprehensive utilization of the diverse multidimensional irregularities inherent in position embedding.

\textbf{Search space}. We construct an extensive search space that captures both types of non-uniformity. This is detailed in Table~\ref{tbl:searchspace}. More precisely, the search framework accommodates the selection of distinct rescaling factors for each dimension within the RoPE embedding. To streamline the structure of the search space, the parameters $\lambda_i$ and $\hat{n}$ are directly explored rather than the indicator function $\mathbb{I}(\hat{\lambda}_{i},\hat{n})$, where the relationship $\hat{\lambda}_{i}=1/\lambda_i$ holds. As detailed in Table~\ref{tbl:searchspace}, the range for $\lambda_i$ spans from a lower bound of 1.0 (corresponding to straightforward extrapolation) to an upper bound of $s\times1.25$ (which represents a broader interpolation scope than PI), using a granularity of 0.01 for the step size, with $s$ denoting the targeted ratio for context window enlargement.

The parameter $\hat{n}$ determines how many of the initial token positions preserve the original RoPE embedding, foregoing position interpolation. In practice, we enable $\hat{n}$ to be selected from the set \{0, 1, 2, 4, 8, 12, 16, 20, 24, 28, 32, 64, 128, 256\}. If $\hat{n}=0$, then the rescale factors identified through search are applied to every token position.

\textbf{Evolution-based search}. The search space detailed in Table~\ref{tbl:searchspace} encompasses a vast number of options for positional interpolation, making comprehensive exploration highly challenging. For instance, when utilizing an extension of $s=4\times$, the possible configurations amount to $400^{128/2}\times14=4\times10^{167}$. This search space grows exponentially as the extension ratio increases. To manage this complexity, we employ evolution search~\cite{guo2020single} and integrate two optimization strategies that substantially enhance search efficiency. The entire search workflow is summarized in Algorithm~\ref{alg:search}.

\textit{Optimized initial population generation}. Rather than generating all $P$ rescale factors at random, we explicitly incorporate the three RoPE rescale factors for PI, NTK, and YaRN as members of the initial population. The other $P$-3 individuals are created by applying random mutations to these three rescale factors, each with a probability of $p$.

\textit{Monotonically non-decreasing constraint}. Upon creating the initial population, LLM perplexity is evaluated for each candidate. This involves applying their specific RoPE rescale factors within the target LLM and calculating perplexity on input $\mathbf{X}$. The individuals with the lowest perplexity—top-k—are selected as parents for the evolutionary step. Nonetheless, due to the immense search space, random mutation and crossover may direct the search toward suboptimal regions, resulting in excessive and redundant perplexity computations. This inefficiency becomes more pronounced when $L'$ is large, as each perplexity evaluation demands considerable inference time.

To mitigate this issue, we enforce a constraint that the sampled RoPE rescaling factors must be monotonically non-decreasing, such that $\lambda_i\le \lambda_{i+1}$ holds for all $i$. Only those RoPE configurations fulfilling this condition are considered during LLM perplexity assessments, which substantially narrows the search space. More precisely, we stipulate that $\lambda_i$ follows a monotonic increase with respect to the RoPE dimension index ($i$ = 0,...,63). This requirement for dimensional monotonicity is motivated by findings in NTK theory~\cite{jacot2018neural,tancik2020fourier,ntk}: higher frequency (lower dimension) components benefit from reduced interpolation (hence smaller $\lambda_i$), whereas lower frequency (higher dimension) components permit greater interpolation (thus, larger $\lambda_i$).

\begin{algorithm}[t]
	\small
	\caption{The search algorithm}
	\textbf{Input:} target LLM, input samples $\mathbf{X}$, population size $P$, mutation size $N_1$, crossover size $N_2$, maximum iterations $\mathcal{T}$, mutation probability $p$\\

	\begin{algorithmic}[1]
		\label{alg:search}
		\STATE Top-k $\leftarrow \phi$;
		\STATE $\text{P}_0$ $\leftarrow$ \textit{Initialize\_population\_with\_optimization}($P$, $\mathbf{X}$, $p$); 
		\FOR{$i$ from $1$ to $\mathcal{T}$}
		\STATE \textit{Compute\_perplexity}(LLM, $\text{P}_{i-1}$, $\mathbf{X}$);
		\STATE Top-k $\leftarrow$ \textit{Update\_Topk}(Top-k, $\text{P}_{i-1}$);
		\STATE $\text{P}_{mutation}$ $\leftarrow$ \textit{Mutation\_with\_mono\_constraint}(Top-k, $N_1$, $p$);
		\STATE $\text{P}_{crossover}$ $\leftarrow$ \textit{Crossover\_with\_mono\_constraint}(Top-k, $N_2$);
		\STATE $\text{P}_i$ $\leftarrow$ $\text{P}_{mutation}$ $\cup$ $\text{P}_{crossover}$ $\cup$ Top-k;
		\ENDFOR
		\STATE Return the element within Top-k exhibiting the minimal perplexity;
	\end{algorithmic}
\end{algorithm}

\textbf{8$\times$ extension without fine-tuning}. Through the application of evolutionary search, our approach successfully discovers non-uniform RoPE rescale factors, maintaining essential dimensions and positions to reduce the information degradation caused by interpolation. As illustrated in Fig.\ref{fig:non-ft}, this strategy allows for expanding LLaMA2’s context window from 4k up to 32k tokens without requiring fine-tuning. In comparison, alternative techniques like PI, as well as non-uniform NTK and YaRN, result in a marked increase in perplexity once the context window is extended beyond 2$\times$.

\subsection{Extending LLM Context Window to 2048K}
\label{sec:extendto2048k}

\textbf{Progressive extension to 2048k}. Here, we present our approach for expanding the context window of pre-trained LLMs from the standard 4k tokens up to beyond 2048k tokens. Our results indicate that employing non-uniform positional interpolation enables an 8$\times$ context extension without requiring fine-tuning. To achieve substantially larger context windows (such as a 512$\times$ increase), however, fine-tuning becomes essential. One possible strategy involves identifying optimal RoPE rescaling factors tailored to the target 2048k length prior to fine-tuning. Nevertheless, this technique presents significant obstacles due to the immense computational costs involved in training. Furthermore, according to our observations, effectively fine-tuning LLMs with such a large extension ratio proves to be difficult (refer to Appendix).

Fortunately, LongRoPE demonstrates efficacy in enhancing both the base and further extended versions of the LLM. Consequently, we propose an efficient and gradual approach that enables reaching the desired 2048k target using only 1k fine-tuning iterations, all conducted within a maximum training length of 256k.

$\Diamond$ \textit{LongRoPE search for extending pre-trained LLM context window to 256k.} Using LLaMA2 as a representative case, we perform a search procedure to expand the context window to 128k and 256k tokens. At this phase, the extension corresponds to 32$\times$ and 64$\times$ increases over the original size, respectively.

$\Diamond$ \textit{Fine-tuning to 256k}. Subsequently, we adapt the pre-trained LLM for a 256k context window through fine-tuning. Our procedure begins with 400 fine-tuning steps on LLaMA2, applying RoPE rescaled factors for 128k. Following this stage, we update the RoPE rescaled factors to accommodate 256k at the obtained checkpoint and continue fine-tuning for another 600 steps. This staged approach demonstrates greater efficiency compared to directly fine-tuning for 256k context length.

$\Diamond$ \textit{Scaling fine-tuned extended LLM up to 2048k using LongRoPE search}. In the final step, we conduct an additional search process based on the LLM that has already been fine-tuned to a 256k context length. Through this approach, we achieve an exceptionally large context window of 2048k, all without necessitating any additional fine-tuning. The overall extension ratio reaches $512\times$.

\textbf{Performance within reduced context window}. Expanding the context window to a substantial 2048k length leads to diminished performance inside the initial window size. This phenomenon is a recognized limitation of positional interpolation~\cite{pi}, which compresses positional embeddings for the original window into a limited space in higher dimensions, ultimately degrading the model’s output quality. Specifically, employing a 512$\times$ expansion causes positions inside the former 4k window to be densely packed.

To address this issue, we conduct an additional evolutionary search on the extended LLM to fine-tune the RoPE rescale parameters specifically for shorter context windows (such as 4k and 8k tokens). For these shorter ranges, we constrain the upper limit for $\lambda$ in the search space, since only minimal positional interpolation is necessary. At inference time, the LLM adaptively modifies the appropriate RoPE rescale factors as needed.

 \section{Experiments}

\label{sec:exp}
\subsection{Setup}
\textbf{Evaluation Tasks and Models}. Our study utilizes LongRoPE with LLaMA2-7B and Mistral-7B, assessing their capabilities through three main evaluation avenues: (1) examining the perplexity of long-context LLMs when processing lengthy documents; (2) conducting the Passkey retrieval test, which evaluates how well a model can extract a straightforward passkey embedded among a large volume of unrelated text; and (3) applying conventional LLM benchmarks within a restricted context window of 4096 tokens.

\textbf{Fine-tuning}. For LLaMA2, the training is conducted with a learning rate of 2e-5 that linearly decays, utilizing a global batch size of 32. The initial phase involves 400 steps of fine-tuning on the Redpajama~\cite{redpajama} dataset, divided into 128k token sequences, each framed by BOS and EOS markers. Subsequently, starting from the obtained checkpoint, a further 600 steps are carried out to extend the context length to 256k tokens. The 128k token context is fine-tuned on 8 A100 GPUs using the distributed training approach described in~\cite{cube}; scaling up to 256k tokens requires 16 A100 GPUs. For Mistral, a fixed learning rate of 1e-6 and a global batch size of 64 are adopted. Fine-tuning for both 128k and 256k models adheres to the procedure outlined in YaRN~\cite{yarn}, involving 400 training steps on Together Computer's Long-Data Collections~\cite{mistraldata}, with sequences of length 16k, using 4 A100 GPUs.

\textbf{Search}. For target window sizes up to 256k, the configuration used is: $P$=64, $N_1$=$N_2$=16, $p$=0.3, $\mathcal{T}$=40, with the top 32 candidates chosen for mutation and crossover in every iteration. Perplexity is determined from 5 randomly chosen validation examples from the PG19 set, each satisfying at least the target context length. For window sizes exceeding 512k, the sizes for the population, mutation, and crossover are reduced by half. In this scenario, perplexity assessment employs 3 randomly selected samples from the Pile-Books3~\cite{pile} validation set.

\textbf{Baselines}. To achieve a context window size of 2048k, we conducted fine-tuning on models with initial context windows of 128k and 256k. Accordingly, we refer to these as LongRoPE-2048k (ft=128k) for LLaMA2 and LongRoPE-2048k (ft=256k) for Mistral. Our evaluation includes a comparison of all four models against leading baselines in context window extension, particularly focusing on open-source LLMs that have undergone positional interpolation fine-tuning via PI, NTK, and YaRN methods. Among these are Together-32k~\cite{together}, Code LLaMA~\cite{codellama}, LongLoRA-full-FT-100k~\cite{longlora}, YaRN-LLaMA, and YaRN-Mistral~\cite{yarn}.

\subsection{Main Results}

\label{sec:mainresult}
\textbf{Modeling long sequences up to 256k tokens}. Our initial analysis assesses our approach against leading-edge extended LLMs on tasks constrained to a 256k token window. To highlight the robustness of our model across diverse settings, we conduct experiments using the test splits from two datasets: Proof-pile~\cite{pg19} and PG19~\cite{pile}. We report perplexity measured across multiple context lengths, employing a sliding window of size 256 for evaluation. For PG19, evaluation is performed on all 100 documents in the test set. In the case of Proof-pile, we follow the sample selection protocol of YaRN~\cite{yarn}, drawing 10 samples at random, each comprising at least 128k tokens.

Table~\ref{tbl:proof-pile} and Table~\ref{tbl:pg19} present a perplexity comparison between LLaMA2 and Mistral models, both of which have been augmented using various interpolation techniques, across the Proof-pile and PG19 datasets. From these results, two main findings emerge: \textbf{(1)} our expanded models consistently demonstrate a reduction in perplexity as the evaluation length increases from 4k to 256k, indicating their enhanced ability to utilize longer contexts. \textbf{(2)} Notably, despite operating with context windows up to 16$\times$ larger—a scenario often associated with degraded performance for short sequences—our LongRoPE-2048k models still achieve lower perplexity compared to leading baselines when evaluated at a 256k context length.

\textbf{Language modeling for sequences exceeding 2000k tokens}. In order to assess performance on exceptionally lengthy documents, we utilize the Books3~\cite{pile} dataset. To streamline evaluation, we draw a random sample of 20 books with lengths over 2048k, applying a 256k sliding window for processing.

As illustrated in Table~\ref{tbl:books3}, LongRoPE is able to effectively increase the context window of both LLaMA2-7B and Mistral-7B up to 2048k tokens, while maintaining perplexity levels that are on par with or surpass those of baseline methods across the 8k-128k range. Clear differences emerge between the models when using the 2048k window. Mistral displays stronger results than baselines at shorter context lengths, yet its perplexity rises above 7 for sequences exceeding 256k tokens. For LLaMA2, the observed perplexity trends agree with expectations: performance consistently improves—perplexity drops—as the context grows, although slight increases are noted at the 1024k and 2048k settings. Additionally, for LLaMA2, LongRoPE-2048k achieves superior results when fine-tuned at 256k rather than 128k, which can be attributed to its use of a smaller secondary extension ratio (8$\times$ instead of 16$\times$). Conversely, Mistral achieves better results with a fine-tuning window of 128k. This primarily stems from following YaRN's procedure during Mistral's 128k and 256k fine-tuning, where a relatively short training sequence length of 16k is used; this choice impacts Mistral's capacity for further context expansion post fine-tuning.

\textbf{Passkey retrieval}. Next, we examine the practical limitations of context window length in generation scenarios. To do so, we utilize the synthetic passkey retrieval task described by ~\cite{passkey}. Here, the model must locate and return a randomly generated five-digit number embedded somewhere in an extended passage. The exact prompt construction is outlined in the appendix. For our experiments, we conduct 10 rounds of the passkey retrieval test, each time positioning the passkey at a random point, chosen uniformly from the available evaluation context.

Fig.~\ref{fig:passkey} presents a comparison of retrieval accuracy against baseline models. The accuracy of previously established approaches quickly falls to zero once the sequence length exceeds 128k. By contrast, even when faced with the demanding task of extracting a passkey from token sequences reaching the million scale, LongRoPE-LLaMA2-2048k (ft=256k) is able to consistently deliver high retrieval accuracy (at least 90\%) across sequence lengths ranging from 4k up to 2048k. Similarly, LongRoPE-Mistral-2048k (ft=128k) maintains perfect (100\%) retrieval accuracy up to 1800k, only experiencing a decrease to 60\% at 2048k—behavior that corresponds with the results shown in Table~\ref{tbl:books3}, where perplexity shows a slight increase at the maximum length.

\textbf{Standard benchmarks under the original context window}. We assess the performance of LongRoPE-2048k models within the original context window by utilizing the Hugging Face Open LLM Leaderboard~\cite{huggingfaceboard} in both zero-shot and few-shot scenarios. Specifically, our evaluation incorporates 25-shot ARC-Challenge~\cite{allenai:arc}, 10-shot HellaSwag~\cite{zellers2019hellaswag}, 5-shot MMLU~\cite{mmlu}, as well as 0-shot TruthfulQA~\cite{lin2021truthfulqa}.

According to Table~\ref{tbl:commonsense}, our models perform similarly to the baseline on the original benchmark, which was created for use with a shorter context window. Notably, our models surpass the original Mistral on TruthfulQA by a margin of +0.5\%. Although LongRoPE-LLaMA2-2048k, which was fine-tuned at 256k, experiences a modest decline in performance, its results largely stay within acceptable limits across the majority of evaluated tasks.

\subsection{Ablation Results}

\textbf{Effectiveness of the second positional interpolation}. Within our progressive extension framework, we apply our search-based method to implement a second round of non-uniform positional interpolation on large language models that have already been fine-tuned and extended. To assess the impact of this approach, we carry out evaluations using our fine-tuned LLaMA2-256k model. This model is further expanded to 512k, 1024k, and 2048k positions utilizing PI and YaRN for comparison. Results presented in Table~\ref{tbl:secondsearch} indicate that our non-uniform positional interpolation maintains stable perplexity outcomes across expansions. Conversely, perplexity metrics for both PI and YaRN deteriorate rapidly as the extension size increases.

\textbf{Effectiveness of recovery at shorter context lengths.} To address the issue of degraded performance at reduced context lengths, we reoptimize the RoPE factors for LongRoPE-2048k utilizing our search methodology. In this process, we deliberately lower the upper bound on scaling factors, thereby minimizing interpolation when processing shorter 4k and 8k contexts. Table~\ref{tbl:shortperformance} reports the perplexity of LongRoPE-LLaMA2-2048k on Proof-pile for 4k and 8k tokens, as well as its mean accuracy across multiple LLM benchmarks. These findings indicate notable gains in performance for short context scenarios.

\textbf{Analysis on the two forms of non-uniformities}. In this section, we conduct ablation studies to individually assess the impact of the two non-uniformities on performance. Two sets of experiments are performed: (i) we extend LLaMA2-7B to shorter sequence lengths of 16k and 32k using several approaches—namely, PI, optimizing only the RoPE dimension, and jointly optimizing both non-uniformities; (ii) we further extend our 256k fine-tuned LLaMA2 model up to 2048k by following a similar methodology. Perplexity scores are measured without additional fine-tuning. As illustrated in Table~\ref{tbl:searchdimension}, optimizing for non-uniformity in RoPE dimension yields notably lower perplexity than the linear interpolation employed by PI. Adjusting non-uniformity for token positions leads to improved results at 16k and 32k, although this benefit is not seen at 2048k, likely owing to the much greater sequence length. Solely keeping the original tokens, without performing any interpolation, does not yield advantages and is reserved for future investigation.

\section{Related Works}

Besides position interpolation methods, this section reviews alternative approaches proposed in related literature.

\textbf{Retrieval-based} strategies utilize an external memory mechanism to store information from earlier contexts, and employ retrieval systems to access pertinent documents during inference~\cite{longllama,wang2023augmenting,borgeaud2022improving}. These frameworks often require substantial alterations to the architectures of LLMs. In comparison, our approach maintains a lightweight design, introducing only minimal changes to the positional embeddings. Furthermore, our method extends applicability to a broader range of long-context tasks beyond retrieval, including areas like long document summarization and few-shot learning.

\textbf{Attention-based context window extensions}. In addition to interpolating positional embeddings, several studies have managed to increase the input context length of LLMs—while retaining the original context window—by modifying attention mechanisms~\cite{han2023lminfinite, streamingllm,parallelwindow}. Central to these approaches is the use of innovative attention masking strategies aimed at alleviating the issue of attention explosion introduced by new position tokens. Such strategies serve as a complement to positional interpolation techniques.

\textbf{Fine-tuning based approaches} investigate strategies for adapting pre-trained LLMs to accommodate extended contexts through adjustments in position embeddings. Research efforts such as Code LLaMA~\cite{codellama}, LLaMA2 Long~\cite{xiong2023effective}, and ScaledRoPE~\cite{liu2023scaling} adopt a substantially larger base value in RoPE, subsequently fine-tuning the model for the desired sequence length. The method presented here allows for adaptable adjustment to a range of target lengths and demonstrates scalability up to lengths exceeding 2M. In light of the considerable computational cost associated with fine-tuning for lengthy contexts (greater than 128k), recent solutions like LongLoRA~\cite{longlora} and PoSE~\cite{zhu2023pose} have emerged to alleviate such requirements. Our approach remains compatible and complementary to these resource-efficient fine-tuning methods.

\section{Conclusion}

In this paper, we introduce LongRoPE, a technique that significantly increases the context window of LLMs to an extraordinary 2048k, all while preserving their performance within the original, smaller context range. Our approach leverages two distinct types of non-uniformity present in the RoPE positional embedding, employing an efficient evolutionary search process. This strategy delivers two key advantages: it supplies robust initialization for fine-tuning and allows for an 8$\times$ expansion of the context window without necessitating fine-tuning. Furthermore, we put forward a progressive extension framework that utilizes LLMs fine-tuned at a 256k context length to systematically achieve a 2048k window size without additional fine-tuning. Our comprehensive experiments demonstrate the efficacy of LongRoPE. We anticipate that the LongRoPE-2048k models will unlock novel long-context tasks and stimulate further exploration in this area.

\section*{Broader Impacts} The research described in this paper aims to contribute to progress in the field of Machine Learning. While the study may have a range of societal implications, we do not identify any particular issues that require special attention at this time.

	\balance

\end{document}